\section*{v1.9.0 Readout Stabilization and Proof Memory Core}
\addcontentsline{toc}{section}{v1.9.0 Readout Stabilization and Proof Memory Core}
\begin{quote}
\textbf{Normalization note.} This block is retained as a historical precursor of the normalized readout architecture. Its proof-bearing content is re-exported later through the grouped-layer normalization, especially the readout/collapse phases and the semantic--tensor transport discipline. It should therefore be read as historical proof memory rather than as an independent new foundation.
\end{quote}

This release adds a further stabilization layer to the inherited closure, probe,
projection, and appendix architecture. Its purpose is to make the final readout
path more disciplined: not only must admissible content survive transport, it
must also remain recoverable after compression. In this sense v1.9.0 pushes
the document from a transport-stable monolith toward a proof-memory stable one.
A statement is not treated as mature merely because it can be compressed into a
lemma, table row, or projection line; it must also remain reconstructible as the
same closure-supported content when the compression is unfolded again.

\subsection*{1. Readout stabilization principle}
The main strengthening of v1.9.0 is that every admissible readout must be
stable under re-expansion. A readout is the visible endpoint of a structural
chain, but it remains acceptable only if its origin can be reopened without
changing the admissibility class. Thus the document now treats readout as a
\emph{controlled fixed point} of compression and reconstruction:
\[
\begin{aligned}
&\text{local structure} \Longrightarrow \text{HC} \Longrightarrow \text{closure}\\
&\Longrightarrow \text{probe filtration} \Longrightarrow \text{table/tool compression}\\
&\Longrightarrow \text{readout} \Longrightarrow \text{re-expansion to the same admissibility class}.
\end{aligned}
\]

\paragraph{Stabilization Proposition A (readout fixed-point discipline).}
If a statement is compressed into a readout form by admissibility-preserving
transports and can be re-expanded to the same closure-licensed content, then the
readout is structurally stable. If re-expansion changes the admissibility class,
the readout was only heuristic.

\paragraph{Proof sketch.}
The admissibility source is fixed upstream by the ground condition, HC, and
closure. Compression is legitimate only if it does not sever that source. Hence
re-expansion must recover the same structural ancestry rather than a merely
similar narrative.

\subsection*{2. Tool-to-lemma synchronization}
Previous versions upgraded tools and tables into proof interfaces. v1.9.0
adds a stricter synchronization rule: every reusable tool should point to at
least one lemma-level justification, and every repeated lemma pattern should be
compressible into a named tool or interface move once stabilized. The goal is to
prevent the appendix from splitting into disconnected ``tricks'' on one side and
formal arguments on the other.

\paragraph{Stabilization Proposition B (tool/lemma synchronization).}
A reusable tool is mature only when it is synchronized with a lemma pattern that
licenses its use, and a recurring lemma pattern is fully normalized only when it
admits a stable compressed tool form.

\paragraph{Proof sketch.}
Tools without lemma ancestry become opaque shortcuts; lemmas without compressed
interface forms bloat the running line unnecessarily. Synchronization preserves
both transparency and speed while keeping the proof architecture auditable.

\subsection*{3. Probe-to-readout compatibility}
The probe channel now receives a further role. It is not only a gate before
projection but also a compatibility filter for later readouts. Any sector that
passes through probe filtration should emerge with enough structural trace to be
placed into a table, a compressed lemma, or a projection sentence without losing
its closure ancestry.

\paragraph{Stabilization Proposition C (probe-filtered modes carry readout ancestry).}
If a mode survives the probe filtration under the standing orthogonality and HC
rules, then its later readout must retain a backward trace to the probe channel
that selected it. Readout therefore cannot erase probe ancestry without lowering
its status to heuristic interpretation.

\subsection*{4. Appendix as proof memory lattice}
The appendix is sharpened once more: it is now treated as a \emph{proof memory
lattice}. Its entries are not merely stored results but nodes from which one can
move upward to the local closure chain, sideways to equivalent interface forms,
and downward to bounded projection language. This lattice viewpoint organizes the
appendix as a recoverability structure rather than a list of leftovers.

\paragraph{Stabilization Principle (memory-lattice admissibility).}
An appendix entry is fully admissible only if it supports three motions:
\begin{enumerate}[leftmargin=2em,label=(\alph*)]
  \item upward recovery to its closure/lemma ancestry,
  \item lateral transport to equivalent tool or table interfaces, and
  \item downward use in a bounded readout or projection context.
\end{enumerate}
If one of these motions fails, the entry must be marked as provisional,
heuristic, or not yet fully normalized.

\subsection*{5. Operational consequence for later versions}
With v1.9.0 the monolith receives a clearer maturity criterion for future
expansions. New derivations should not merely add content; they should improve
recoverability. In practice this means that later mathematical deepening should
prefer additions that tighten one of the following links: closure to probe,
probe to table, table to lemma, lemma to tool, or readout back to structural
ancestry. This keeps the entire project aligned with the structure-first rule
while making the growing proof apparatus easier to audit and reuse.

\paragraph{Operational rule for subsequent releases.}
Whenever a new argument is introduced, one should ask not only whether the step
is locally correct, but also whether it improves the global recoverability of
compressed proof content. The strongest additions are those that make the same
statement simultaneously easier to compress, easier to re-expand, and easier to
locate inside the appendix memory lattice.



